\section{TIẾN TRÌNH LỊCH SỬ HÌNH THÀNH VÀ PHÁT TRIỂN HỆ THỐNG GIÁO DỤC NGHỀ NGHIỆP VIỆT NAM TỪ NĂM 1945 ĐẾN NAY}

Tiến trình hình thành và phát triển của hệ thống giáo dục nghề nghiệp Việt Nam từ mùa thu Cách mạng Tháng Tám năm 1945 đến nay là một chặng đường lịch sử hào hùng nhưng cũng đầy thử thách, phản ánh trực tiếp sự thích ứng của hệ thống giáo dục và đào tạo trước những đòi hỏi của từng giai đoạn cách mạng. Để làm sáng tỏ quy luật vận động của hệ thống, chương này nghiên cứu tiến trình lịch sử qua sáu giai đoạn phân kỳ cơ bản, gắn liền với sự biến đổi của bối cảnh kinh tế -- xã hội, sự điều chỉnh của thể chế quản lý nhà nước và việc ban hành các văn bản quy phạm pháp luật nền tảng.

\begin{figure}[htbp]
\centering
\resizebox{0.95\linewidth}{!}{%
\begin{tikzpicture}[
    node distance = 0.85cm and 0cm,
    box/.style = {
        rectangle, 
        rounded corners=4pt, 
        draw=blue!70!black, 
        fill=blue!4, 
        line width=0.8pt, 
        text width=14.2cm, 
        align=center,
        inner sep=6pt,
        font=\small
    },
    arrow/.style = {
        -{Stealth[scale=1.1]}, 
        line width=1.1pt, 
        draw=blue!60!black
    }
]

% Node 1
\node (n1) [box] {
    \textbf{Giai đoạn 1945 -- 1954: Kháng chiến kiến quốc và đặt nền móng}\\
    \textit{Cơ quan chủ quản:} Nha Kỹ thuật vụ (Bộ Quốc gia Giáo dục, theo Sắc lệnh số 119/SL)\\
    \textit{Mô hình:} Đào tạo công nhân kỹ thuật quốc phòng, phục vụ công xưởng chiến khu
};

% Node 2
\node (n2) [box, below=of n1] {
    \textbf{Giai đoạn 1954 -- 1975: Xây dựng chủ nghĩa xã hội ở miền Bắc và chi viện miền Nam}\\
    \textit{Cơ quan chủ quản:} Vụ Quản lý nhân công $\rightarrow$ Tổng cục Đào tạo công nhân kỹ thuật (1969)\\
    \textit{Mô hình:} Đào tạo theo ngành công nghiệp nặng; hợp tác với Liên Xô và Đông Âu
};

% Node 3
\node (n3) [box, below=of n2] {
    \textbf{Giai đoạn 1975 -- 1986: Thống nhất đất nước và kinh tế kế hoạch hóa bao cấp}\\
    \textit{Cơ quan chủ quản:} Tổng cục Dạy nghề trực thuộc Hội đồng Chính phủ (Nghị định số 240/CP)\\
    \textit{Mô hình:} Đào tạo theo chỉ tiêu pháp lệnh; Nhà nước bao cấp và phân bổ việc làm
};

% Node 4
\node (n4) [box, below=of n3] {
    \textbf{Giai đoạn 1986 -- 2006: Bước ngoặt Đổi mới và chuyển giao thể chế}\\
    \textit{Cơ quan chủ quản:} Bộ Giáo dục và Đào tạo $\rightarrow$ Tổng cục Dạy nghề thuộc Bộ Lao động (1998)\\
    \textit{Mô hình:} Đa dạng hóa loại hình cơ sở dạy nghề thích ứng kinh tế thị trường (Quyết định 253)
};

% Node 5
\node (n5) [box, below=of n4] {
    \textbf{Giai đoạn 2006 -- 2014: Thể chế hóa chuyên biệt theo Luật Dạy nghề 2006}\\
    \textit{Cơ quan chủ quản:} Tổng cục Dạy nghề (Bộ Lao động -- Thương binh và Xã hội)\\
    \textit{Mô hình:} Ba cấp trình độ (Sơ cấp nghề, Trung cấp nghề, Cao đẳng nghề); Đề án 1956
};

% Node 6
\node (n6) [box, below=of n5] {
    \textbf{Giai đoạn 2014 -- nay: Hợp nhất toàn diện theo Luật Giáo dục nghề nghiệp 2014}\\
    \textit{Cơ quan chủ quản:} Tổng cục Giáo dục nghề nghiệp (Bộ Lao động -- Thương binh và Xã hội)\\
    \textit{Mô hình:} Hệ thống giáo dục nghề nghiệp mở, liên thông; Khung trình độ quốc gia (VQF)
};

% Lines with arrows
\draw [arrow] (n1) -- (n2);
\draw [arrow] (n2) -- (n3);
\draw [arrow] (n3) -- (n4);
\draw [arrow] (n4) -- (n5);
\draw [arrow] (n5) -- (n6);

\end{tikzpicture}%
}
\caption{Tiến trình lịch sử và sự chuyển biến thể chế quản lý hệ thống giáo dục nghề nghiệp Việt Nam (1945 -- nay)}
\label{fig:so_do_tien_trinh}
\end{figure}


\subsection{Giai đoạn 1945 -- 1954: Đào tạo công nhân kỹ thuật trong kháng chiến chống thực dân Pháp và kiến quốc}

Cách mạng Tháng Tám năm 1945 thành công đã khai sinh ra nước Việt Nam Dân chủ Cộng hòa, mở ra một kỷ nguyên mới cho dân tộc. Tuy nhiên, chính quyền non trẻ ngay lập tức phải đối mặt với muôn vàn khó khăn hiểm nghèo, thường được ví như tình thế ngàn cân treo sợi tóc. Đất nước phải đồng thời đương đầu với ba thứ giặc nguy hiểm là giặc đói, giặc dốt và giặc ngoại xâm. 

Về phương diện kinh tế và kỹ thuật, di sản mà chế độ thực dân để lại là một nền kinh tế lạc hậu cùng một cơ sở hạ tầng công nghiệp què quặt. Tỷ lệ mù chữ trong nhân dân lên tới hơn chín mươi lăm phần trăm. Đội ngũ thợ kỹ thuật người Việt Nam lúc bấy giờ vô cùng mỏng manh, chủ yếu làm việc tại một số cơ sở hỏa xa, xưởng sửa chữa tàu thủy nhỏ hoặc các hầm mỏ khai thác than, trình độ chuyên môn kỹ thuật phần lớn mang tính kinh nghiệm thực hành hạn hẹp, thiếu hụt nghiêm trọng kiến thức lý thuyết cơ bản.

Trước tình hình thực dân Pháp quay lại xâm lược, nhiệm vụ hàng đầu của Nhà nước cách mạng là phải nhanh chóng xây dựng lực lượng vũ trang nhân dân và tự lực tự cường sản xuất vũ khí, khí tài quân sự. Ngày 09 tháng 7 năm 1946, Chủ tịch Chính phủ Việt Nam Dân chủ Cộng hòa đã ký Sắc lệnh số 119/SL quy định về tổ chức của Bộ Quốc gia Giáo dục. Trong cơ cấu tổ chức bộ máy này, \textit{Nha Kỹ thuật vụ} lần đầu tiên được thành lập với trọng trách quản lý và chỉ đạo công tác đào tạo cán bộ kỹ thuật, công nhân kỹ thuật phục vụ cho nhu cầu kháng chiến và kiến quốc. Đây được coi là mốc pháp lý đầu tiên khẳng định sự hình thành của cơ quan quản lý chuyên trách về giáo dục kỹ thuật và dạy nghề trong chính quyền nhân dân.

Khi cuộc kháng chiến toàn quốc bùng nổ vào tháng 12 năm 1946, các cơ sở sản xuất và trường lớp kỹ thuật buộc phải triệt để sơ tán lên các chiến khu bí mật tại Việt Bắc, Liên khu 4 và Liên khu 5. Dưới sự bảo bọc của đồng bào, các xưởng quân giới và các lớp đào tạo thợ kỹ thuật di động được hình thành ngay giữa lòng rừng núi. Tiêu biểu là các cơ sở trường lớp gắn liền với Cục Quân giới do Giáo sư, Viện sĩ Trần Đại Nghĩa lãnh đạo. 

Mô hình đào tạo trong giai đoạn này mang tính chất đặc thù thời chiến rất rõ nét: kết hợp chặt chẽ việc dạy lý thuyết ngắn gọn với thực hành ngay trên dây chuyền sản xuất vũ khí; phương châm hành động là người biết nhiều dạy người biết ít, vừa chiến đấu, vừa sản xuất, vừa học tập. Mặc dù trang thiết bị vô cùng thô sơ, chủ yếu là các máy tiện, máy phay cũ kỹ được tháo dỡ và vận chuyển gian khổ từ các thành phố lớn lên chiến khu, đội ngũ công nhân kỹ thuật được tôi luyện trong khói lửa đã làm chủ được các công nghệ khó, chế tạo thành công súng bazooka, súng không giật SKZ, mìn gạt và lựu đạn phục vụ hiệu quả cho các chiến dịch quân sự lớn.

Bước sang năm 1950, khi cuộc kháng chiến bước sang giai đoạn phản công mạnh mẽ, Trung ương Đảng và Chính phủ đã quyết định tiến hành cuộc Cải cách giáo dục năm 1950. Ngày 11 tháng 7 năm 1950, Chủ tịch nước ban hành Sắc lệnh số 121/SL và Sắc lệnh số 122/SL xác lập cấu trúc nền giáo dục mới của nước Việt Nam Dân chủ Cộng hòa. 

Cuộc cải cách giáo dục này có ý nghĩa lịch sử sâu sắc khi lần đầu tiên xác định hệ thống giáo dục quốc dân mới gồm hai bộ phận hữu cơ gắn bó mật thiết với nhau: giáo dục phổ thông và giáo dục chuyên nghiệp (bao gồm các trường trung học kỹ thuật và các trường thực hành nghề). Các trường thực hành kỹ nghệ thời kỳ này đã được củng cố lại về chương trình, đưa thêm kiến thức khoa học thường thức vào giảng dạy, chuẩn bị sẵn sàng một lực lượng công nhân có kỹ năng nhằm phục vụ cho công cuộc tiếp quản các đô thị lớn sau chiến thắng lịch sử Điện Biên Phủ năm 1954.

\subsection{Giai đoạn 1954 -- 1975: Đào tạo công nhân kỹ thuật phục vụ xây dựng chủ nghĩa xã hội ở miền Bắc và chi viện miền Nam}

Sau Hiệp định Giơ-ne-vơ năm 1954, đất nước tạm thời bị chia cắt làm hai miền với hai chế độ chính trị khác nhau. Miền Bắc bước vào thời kỳ quá độ lên chủ nghĩa xã hội, thực hiện khôi phục kinh tế, cải tạo xã hội chủ nghĩa và triển khai Kế hoạch năm năm lần thứ nhất (1961 -- 1965) với mục tiêu xây dựng bước đầu cơ sở vật chất kỹ thuật của chủ nghĩa xã hội. 

Nhu cầu kiến thiết đất nước đòi hỏi phải xây dựng hàng loạt khu công nghiệp tập trung như khu công nghiệp gang thép Thái Nguyên, khu công nghiệp dệt Nam Định, khu công nghiệp hóa chất Việt Trì, cùng hệ thống các nhà máy cơ khí, điện lực tại Hà Nội và Hải Phòng. Nhu cầu nhân lực kỹ thuật trực tiếp cho các ngành công nghiệp nặng, công nghiệp nhẹ, giao thông vận tải và xây dựng tăng lên theo cấp số nhân, đặt công tác đào tạo công nhân kỹ thuật vào vị trí mặt trận trọng yếu.

Trong những năm đầu sau hòa bình lập lại (từ năm 1955 đến năm 1969), công tác quản lý và phân phối nhân lực lao động kỹ thuật được giao cho Vụ Quản lý nhân công thuộc Bộ Lao động đảm nhiệm. Tuy nhiên, trước sự phát triển vượt bậc của mạng lưới các xí nghiệp quốc doanh và đòi hỏi chuẩn hóa tay nghề ngày càng khắt khe, mô hình quản lý cấp vụ tỏ ra không còn tương xứng với quy mô nhiệm vụ.

Bước ngoặt thể chế mang tính quyết định của giai đoạn này được đánh dấu bằng việc Hội đồng Chính phủ ban hành Nghị định số 200/CP ngày 09 tháng 10 năm 1969 về việc thành lập \textbf{Tổng cục Đào tạo công nhân kỹ thuật} trực thuộc Bộ Lao động. Đây là lần đầu tiên trong lịch sử, một cơ quan quản lý nhà nước tầm cỡ tổng cục được thiết lập riêng biệt để chuyên trách chỉ đạo, quy hoạch và tổ chức toàn diện công tác đào tạo công nhân kỹ thuật trên quy mô toàn miền Bắc. 

Ngay sau đó, ngày 30 tháng 11 năm 1971, Hội đồng Chính phủ tiếp tục ban hành Chỉ thị số 196-CP về công tác đào tạo công nhân kỹ thuật, chỉ rõ yêu cầu các bộ, ngành và địa phương phải xây dựng quy hoạch đào tạo nghề dài hạn gắn chặt với chiến lược phát triển kinh tế quốc dân.

Mô hình đào tạo công nhân kỹ thuật trong giai đoạn 1954 -- 1975 mang những đặc trưng vô cùng độc đáo và chặt chẽ:

Trước hết, đó là sự hình thành mạng lưới trường đào tạo công nhân kỹ thuật trực thuộc các bộ chuyên ngành. Các bộ như Bộ Cơ khí và Luyện kim, Bộ Xây dựng, Bộ Giao thông vận tải, Bộ Điện và Than, Bộ Thủy lợi đều xây dựng các trường công nhân kỹ thuật nòng cốt của ngành mình (tiêu biểu như Trường Công nhân Cơ khí Hà Nội, Trường Công nhân Kỹ thuật Xây dựng, Trường Công nhân Lái xe Đường bộ). Học sinh tốt nghiệp phổ thông được tuyển chọn kỹ lưỡng, được hưởng chế độ bao cấp toàn bộ về tiền ăn, ký túc xá, trang phục bảo hộ và phụ cấp sinh hoạt hàng tháng.

Thứ hai, chuẩn đầu ra tay nghề được thiết lập nghiêm ngặt theo thang bậc thợ bảy bậc (phổ biến là bậc hai trên bảy và bậc ba trên bảy cho công nhân mới ra trường). Chương trình đào tạo cân đối hài hòa giữa lý thuyết chuyên môn nghề với thực hành tại xưởng trường và thực tập trực tiếp trên các công trường, nhà máy.

Thứ ba, sự mở rộng hợp tác quốc tế toàn diện với các nước xã hội chủ nghĩa anh em. Hàng vạn học sinh xuất sắc và thanh niên miền Bắc đã được Nhà nước cử đi đào tạo công nhân kỹ thuật và thực tập sinh nghề tại Liên bang Xô Viết, Cộng hòa Dân chủ Đức, Tiệp Khắc, Ba Lan, Hung-ga-ri và Bun-ga-ri. Lực lượng này khi trở về nước đã trở thành những công nhân lành nghề đầu đàn, những tổ trưởng sản xuất, kỹ sư thực hành chủ chốt, có khả năng làm chủ các dây chuyền thiết bị hiện đại do các nước bạn viện trợ. 

Ngay cả trong những năm tháng đế quốc Mỹ leo thang ném bom đánh phá ác liệt miền Bắc (1965 -- 1972), các trường công nhân kỹ thuật vẫn kiên cường sơ tán vào các vùng nông thôn, vừa duy trì dạy tốt học tốt dưới hầm hào, vừa trực tiếp tham gia chiến đấu bảo vệ giao thông huyết mạch và chi viện hàng vạn thợ cơ khí, thợ lái xe cho chiến trường miền Nam.

\subsection{Giai đoạn 1975 -- 1986: Thống nhất hệ thống dạy nghề và cơ chế kế hoạch hóa tập trung bao cấp}

Mùa xuân năm 1975, cuộc kháng chiến chống Mỹ cứu nước toàn thắng, non sông thu về một mối. Đất nước bước vào kỷ nguyên độc lập, thống nhất và cùng tiến lên chủ nghĩa xã hội. Nhiệm vụ trung tâm của toàn Đảng, toàn dân lúc này là khắc phục hậu quả nặng nề của chiến tranh, khôi phục hạ tầng kinh tế, thống nhất các thể chế văn hóa, xã hội và giáo dục trên phạm vi cả nước.

Tại miền Nam, hệ thống cơ sở dạy nghề trước năm 1975 chủ yếu được tổ chức theo mô hình phương Tây với mạng lưới các trường bá nghệ, trường kỹ thuật công lập và tư thục tại Sài Gòn, Đà Nẵng, Cần Thơ, Vĩnh Long. Cơ quan quản lý nhà nước cách mạng đã khẩn trương tiếp quản nguyên vẹn các cơ sở vật chất này, tiến hành cải tạo nội dung chương trình theo đường lối xã hội chủ nghĩa, đồng thời bổ sung đội ngũ cán bộ giảng dạy từ miền Bắc vào để nhanh chóng ổn định và tái khởi động guồng máy đào tạo nghề.

Để thống nhất một đầu mối quản lý trên phạm vi toàn quốc và nâng cao vị thế của ngành dạy nghề trong công cuộc tái thiết đất nước, ngày 29 tháng 9 năm 1978, Hội đồng Chính phủ đã ban hành Nghị định số 240/CP quy định nhiệm vụ, quyền hạn và tổ chức bộ máy của \textbf{Tổng cục Dạy nghề} trực thuộc Hội đồng Chính phủ. 

Việc nâng cấp Tổng cục Dạy nghề từ một cơ quan trực thuộc Bộ Lao động lên thành một cơ quan trực thuộc trực tiếp Hội đồng Chính phủ (tương đương cấp bộ) là một dấu mốc lịch sử thể hiện tầm nhìn chiến lược của Đảng và Nhà nước đối với giai cấp công nhân. Tổng cục Dạy nghề lúc này được trao thẩm quyền quản lý nhà nước thống nhất đối với toàn bộ các trường công nhân kỹ thuật của các bộ, ngành kinh tế và các trường dạy nghề của các tỉnh, thành phố trong cả nước.

Ngày 11 tháng 01 năm 1979, Bộ Chính trị Ban Chấp hành Trung ương Đảng (khóa IV) đã ban hành Nghị quyết số 14-NQ/TW về cải cách giáo dục. Nghị quyết đã khẳng định mạnh mẽ vị trí của công tác đào tạo công nhân kỹ thuật trong việc tái cơ cấu lực lượng lao động xã hội, nhấn mạnh nguyên lý học đi đôi với hành, giáo dục kết hợp với lao động sản xuất, nhà trường gắn liền với xí nghiệp.

Đặc trưng cốt lõi của hệ thống dạy nghề trong giai đoạn 1975 -- 1986 là vận hành hoàn toàn theo cơ chế kế hoạch hóa tập trung bao cấp:

Nhà nước thông qua Ủy ban Kế hoạch Nhà nước và Tổng cục Dạy nghề để giao chỉ tiêu tuyển sinh mang tính pháp lệnh hàng năm cho từng trường và từng ngành nghề đào tạo. Chỉ tiêu tuyển sinh được cân đối dựa trên nhu cầu biên chế lao động dự kiến của các nhà máy, xí nghiệp quốc doanh và nông lâm trường quốc doanh.

Người học trúng tuyển vào các trường công nhân kỹ thuật được coi là đã bước chân vào hàng ngũ biên chế nhà nước dự bị. Toàn bộ học phí, chi phí ăn ở, quân tư trang và sách vở đều do ngân sách nhà nước bao cấp tuyệt đối. Sau khi tốt nghiệp, một trăm phần trăm học sinh đều được Nhà nước bố trí, điều động công tác về các cơ quan, xí nghiệp mà không phải tự tìm kiếm việc làm trên thị trường.

Mặc dù mô hình này đã hoàn thành xuất sắc sứ mệnh lịch sử trong việc cung cấp một lượng lớn công nhân kỹ thuật cho công cuộc khôi phục đất nước sau chiến tranh, nhưng khi nền kinh tế rơi vào cuộc khủng hoảng trầm trọng đầu thập niên 1980 do cơ chế quan liêu bao cấp kìm hãm, hệ thống dạy nghề cũng dần bộc lộ những hạn chế cố hữu. Việc đào tạo xơ cứng theo chỉ tiêu kế hoạch khiến các trường xa rời thực tế, trang thiết bị dạy học nhanh chóng bị lạc hậu do ngân sách cạn kiệt, học sinh thiếu động lực tự thân phấn đấu nâng cao tay nghề do tâm lý ỷ lại vào sự sắp đặt việc làm của Nhà nước.

\subsection{Giai đoạn 1986 -- 2006: Giáo dục nghề nghiệp trong bước ngoặt Đổi mới và chuyển đổi thể chế quản lý}

Đại hội đại biểu toàn quốc lần thứ VI của Đảng vào tháng 12 năm 1986 đã mở ra bước ngoặt lịch sử vĩ đại cho đất nước khi quyết định thực hiện đường lối Đổi mới toàn diện, chuyển đổi nền kinh tế từ cơ chế kế hoạch hóa tập trung quan liêu bao cấp sang nền kinh tế hàng hóa nhiều thành phần vận hành theo cơ chế thị trường có sự quản lý của Nhà nước theo định hướng xã hội chủ nghĩa.

Sự chuyển đổi mô hình kinh tế đã tác động trực diện và gây ra một cú sốc lớn đối với hệ thống các trường đào tạo công nhân kỹ thuật. Khi Nhà nước xóa bỏ cơ chế bao cấp và chấm dứt chế độ phân công công tác sau tốt nghiệp, các doanh nghiệp quốc doanh phải tự hạch toán kinh doanh, buộc phải cắt giảm lao động dôi dư và ngừng tuyển dụng tự động. Hệ quả là trong những năm cuối thập niên 1980 và đầu thập niên 1990, hệ thống trường dạy nghề rơi vào một cuộc khủng hoảng tuyển sinh trầm trọng chưa từng thấy. Học sinh tốt nghiệp trung học phổ thông và trung học cơ sở không còn mặn mà vào trường nghề; hàng loạt trường nghề của các bộ, ngành và địa phương rơi vào cảnh đìu hiu, không có người học, giáo viên phải bỏ nghề hoặc chuyển công tác, nhiều cơ sở buộc phải sáp nhập hoặc giải thể.

Song hành với cuộc khủng hoảng kinh tế -- xã hội là những đợt sắp xếp, luân chuyển bộ máy quản lý nhà nước về đào tạo nghề diễn ra liên tục:

Năm 1987, thực hiện chủ trương tinh gọn bộ máy của Hội đồng Bộ trưởng, Tổng cục Dạy nghề bị giải thể tư cách cơ quan trực thuộc Chính phủ và sáp nhập vào Bộ Đại học và Trung học chuyên nghiệp để thành lập \textit{Bộ Đại học, Trung học chuyên nghiệp và Dạy nghề}. Cơ quan quản lý chuyên trách về dạy nghề bị thu hẹp quy mô xuống thành Vụ Đào tạo nghề.

Năm 1990, Chính phủ tiếp tục sáp nhập Bộ Giáo dục với Bộ Đại học, Trung học chuyên nghiệp và Dạy nghề để hình thành \textit{Bộ Giáo dục và Đào tạo}. Trong cơ cấu tổ chức mới, mảng đào tạo nghề tiếp tục bị thu gọn thành một bộ phận trực thuộc Vụ Trung học chuyên nghiệp và Dạy nghề. Việc sáp nhập này tuy có ưu điểm là quy về một đầu mối giáo dục quốc dân, nhưng trong thực tế triển khai, nguồn lực và sự ưu tiên của Bộ Giáo dục và Đào tạo chủ yếu dồn cho giáo dục phổ thông và giáo dục đại học, khiến cho mảng đào tạo nghề rơi vào tình trạng bị xem nhẹ, tiếng nói chính sách yếu ớt, cơ chế đào tạo chưa gắn được với thị trường lao động và việc làm.

Nhận thức rõ bất cập sâu sắc này, sau Hội nghị lần thứ hai Ban Chấp hành Trung ương Đảng (khóa VIII) về định hướng chiến lược phát triển giáo dục -- đào tạo và khoa học -- công nghệ, ngày 30 tháng 12 năm 1998, Thủ tướng Chính phủ đã ban hành Quyết định số 253/1998/QĐ-TTg về việc chuyển giao nhiệm vụ quản lý nhà nước về đào tạo nghề từ Bộ Giáo dục và Đào tạo sang Bộ Lao động -- Thương binh và Xã hội. 

Trên cơ sở đó, \textbf{Tổng cục Dạy nghề} chính thức được tái lập trực thuộc Bộ Lao động -- Thương binh và Xã hội. Quyết định lịch sử này đã gắn chặt công tác đào tạo nghề với thị trường lao động, giải quyết việc làm và các chính sách an sinh xã hội, mở ra một thời kỳ hồi sinh và bứt phá mạnh mẽ cho hệ thống dạy nghề Việt Nam.

Dưới sự chỉ đạo của Bộ Lao động -- Thương binh và Xã hội và Tổng cục Dạy nghề, mạng lưới các cơ sở dạy nghề từng bước thích ứng với kinh tế thị trường thông qua chính sách xã hội hóa. Luật Giáo dục năm 1998 và Luật Giáo dục năm 2005 được ban hành đã hoàn thiện khung khổ pháp lý chung, công nhận sự bình đẳng giữa các loại hình trường dạy nghề công lập, dân lập và tư thục. Hàng trăm trung tâm dạy nghề quận, huyện được đầu tư xây dựng mới; các trường dạy nghề được phép thu học phí, tự chủ liên kết với doanh nghiệp để đào tạo theo đơn đặt hàng, khôi phục lại quy mô tuyển sinh và từng bước đáp ứng nhu cầu nhân lực cho các khu công nghiệp, khu chế xuất đang phát triển bùng nổ trên khắp cả nước.

\subsection{Giai đoạn 2006 -- 2014: Thời kỳ thể chế hóa chuyên biệt theo Luật Dạy nghề năm 2006}

Bước vào những năm đầu thế kỷ hai mươi mốt, Việt Nam đẩy mạnh tiến trình hội nhập kinh tế quốc tế đỉnh cao với dấu mốc chính thức gia nhập Tổ chức Thương mại Thế giới (WTO) vào năm 2007. Làn sóng đầu tư trực tiếp nước ngoài (FDI) ồ ạt đổ vào Việt Nam đã tạo ra hàng triệu việc làm công nghiệp mới, nhưng đồng thời cũng đặt ra thách thức gay gắt về chất lượng tay nghề và năng suất lao động. 

Để tạo hành lang pháp lý vững chắc và nâng tầm vị thế của lĩnh vực dạy nghề tương xứng với đòi hỏi của đất nước, ngày 29 tháng 11 năm 2006, Quốc hội khóa XI đã chính thức thông qua \textbf{Luật Dạy nghề số 76/2006/QH11} (có hiệu lực thi hành từ ngày 01 tháng 6 năm 2007). Đây là đạo luật chuyên ngành đầu tiên trong lịch sử lập pháp của nước ta điều chỉnh toàn diện các quan hệ xã hội trong lĩnh vực dạy nghề.

Luật Dạy nghề năm 2006 đã tạo nên những bước đột phá mang tính bản lề cho hệ thống:

Thứ nhất, luật xác lập một cấu trúc hệ thống dạy nghề hoàn chỉnh gồm ba cấp trình độ rõ ràng: \textit{Sơ cấp nghề, Trung cấp nghề và Cao đẳng nghề}. Việc lần đầu tiên danh xưng "Cao đẳng nghề" được luật hóa đã mở ra một nấc thang đào tạo kỹ năng thực hành đỉnh cao, phá vỡ định kiến cho rằng học nghề chỉ là các khóa đào tạo ngắn hạn, bậc thấp.

Thứ hai, mạng lưới các trường cao đẳng nghề và trung cấp nghề phát triển nhanh chóng vượt bậc về cả số lượng và quy mô đào tạo. Hàng loạt trường cao đẳng nghề chất lượng cao được thành lập trên cơ sở nâng cấp từ các trường công nhân kỹ thuật danh tiếng trước đây, được đầu tư cơ sở vật chất, nhà xưởng thực hành hiện đại theo các chuẩn nghề quốc gia và quốc tế.

Thứ ba, sự ra đời và triển khai của các đại dự án, chiến lược mang tầm vóc quốc gia. Ngày 27 tháng 11 năm 2009, Thủ tướng Chính phủ ban hành Quyết định số 1956/QĐ-TTg phê duyệt Đề án "Đào tạo nghề cho lao động nông thôn đến năm 2020" (thường gọi là Đề án 1956). Đây là chính sách an sinh và phát triển nhân lực có quy mô bao phủ lớn nhất lịch sử, hỗ trợ dạy nghề miễn phí cho hàng triệu nông dân, góp phần đẩy nhanh quá trình chuyển dịch cơ cấu lao động từ nông nghiệp sang công nghiệp và dịch vụ. Tiếp đó, ngày 29 tháng 5 năm 2012, Thủ tướng Chính phủ ban hành Quyết định số 630/QĐ-TTg phê duyệt Chiến lược phát triển dạy nghề thời kỳ 2011 -- 2020, xác định mục tiêu đào tạo nghề theo ba cấp độ chuẩn: quốc gia, khu vực ASEAN và quốc tế.

Tuy nhiên, bên cạnh những bước phát triển vượt bậc, giai đoạn 2006 -- 2014 cũng bộc lộ một điểm nghẽn thể chế nghiêm trọng nhất trong lịch sử giáo dục Việt Nam: \textbf{sự tồn tại song song và chia cắt giữa hai hệ thống đào tạo nghề nghiệp}. Cụ thể:

Một nhánh gồm các trường sơ cấp nghề, trung cấp nghề và cao đẳng nghề do Bộ Lao động -- Thương binh và Xã hội (Tổng cục Dạy nghề) quản lý nhà nước theo Luật Dạy nghề năm 2006.

Một nhánh khác gồm các trường trung cấp chuyên nghiệp và các khoa trung cấp, cao đẳng trong các trường đại học do Bộ Giáo dục và Đào tạo quản lý nhà nước theo Luật Giáo dục năm 2005.

Sự chia cắt này đã dẫn đến tình trạng chồng chéo nghiêm trọng về chức năng đào tạo, lãng phí lớn nguồn lực đầu tư công của ngân sách, gây xung đột về văn bằng chứng chỉ và tạo ra sự hoang mang tột độ cho người học cũng như xã hội khi lựa chọn lộ trình học tập. Vấn đề bức thiết đặt ra cho các nhà hoạch định chính sách là phải khẩn trương hợp nhất hai nhánh đào tạo này thành một hệ thống thống nhất.

\subsection{Giai đoạn 2014 đến nay: Hợp nhất toàn diện theo Luật Giáo dục nghề nghiệp năm 2014 và hội nhập quốc tế}

Trước đòi hỏi khách quan của thực tiễn và quán triệt sâu sắc tinh thần Nghị quyết số 29-NQ/TW ngày 04 tháng 11 năm 2013 của Hội nghị lần thứ tám Ban Chấp hành Trung ương Đảng (khóa XI) về đổi mới căn bản, toàn diện giáo dục và đào tạo, ngày 27 tháng 11 năm 2014, Quốc hội khóa XIII đã chính thức thông qua \textbf{Luật Giáo dục nghề nghiệp số 74/2014/QH13} (có hiệu lực thi hành từ ngày 01 tháng 7 năm 2015).

Luật Giáo dục nghề nghiệp năm 2014 được đánh giá là một bước ngoặt lập pháp lịch sử mang tính cách mạng đối với hệ thống nhân lực kỹ thuật Việt Nam:

Điểm đột phá căn bản nhất là đạo luật đã dũng cảm xóa bỏ hoàn toàn ranh giới chia cắt kéo dài hàng thập kỷ giữa "dạy nghề" và "trung cấp chuyên nghiệp". Luật đã chính thức hợp nhất hai hệ thống cũ thành một chỉnh thể thống nhất mang tên \textbf{Giáo dục nghề nghiệp}. Hệ thống mới quy định ba trình độ đào tạo chuẩn hóa: \textit{Sơ cấp, Trung cấp và Cao đẳng}. Theo đó, các trường trung cấp chuyên nghiệp và trường trung cấp nghề được chuyển đổi thành trường trung cấp; các trường cao đẳng chuyên nghiệp và trường cao đẳng nghề được chuyển đổi thành trường cao đẳng. 

Nghị quyết của Quốc hội và sau đó là Nghị định số 14/2017/NĐ-CP của Chính phủ đã khẳng định Bộ Lao động -- Thương binh và Xã hội là cơ quan đầu mối thống nhất thực hiện chức năng quản lý nhà nước về giáo dục nghề nghiệp trong cả nước. Tháng 8 năm 2017, Tổng cục Dạy nghề chính thức được đổi tên thành \textbf{Tổng cục Giáo dục nghề nghiệp}, đánh dấu việc quy tụ toàn bộ mạng lưới các trường cao đẳng, trung cấp (trừ khối ngành sư phạm thuộc Bộ Giáo dục và Đào tạo) về một mối chỉ đạo điều hành xuyên suốt.

Nhằm đưa giáo dục nghề nghiệp Việt Nam hội nhập bình đẳng với khu vực và toàn cầu, ngày 18 tháng 10 năm 2016, Thủ tướng Chính phủ đã ban hành Quyết định số 1982/QĐ-TTg phê duyệt \textit{Khung trình độ quốc gia Việt Nam (VQF)}. Khung trình độ quốc gia gồm tám bậc, trong đó các bậc 1, 2, 3 tương ứng với các trình độ sơ cấp nghề nghiệp; bậc 4 tương ứng với trình độ trung cấp; và bậc 5 tương ứng với trình độ cao đẳng. Việc thiết lập khung trình độ này đã tạo ra ngôn ngữ chung để công nhận kỹ năng nghề với Khung tham chiếu trình độ ASEAN (AQRF) và các hệ thống bằng cấp quốc tế.

Bước sang thập niên 2020, trước những biến động nhanh chóng của thị trường lao động dưới tác động của đại dịch toàn cầu, công nghệ số và làn sóng tự động hóa, Chính phủ tiếp tục ban hành Quyết định số 2239/QĐ-TTg ngày 30 tháng 12 năm 2021 phê duyệt \textit{Chiến lược phát triển giáo dục nghề nghiệp giai đoạn 2021 -- 2030, tầm nhìn đến năm 2045}. Chiến lược đặt ra mục tiêu phát triển hệ thống giáo dục nghề nghiệp mở, linh hoạt, hiện đại, hiệu quả và hội nhập; đến năm 2030 có khoảng một trăm trường chất lượng cao đạt chuẩn khu vực và thế giới.

Đặc biệt, ngày 04 tháng 5 năm 2023, Ban Bí thư Trung ương Đảng đã ban hành \textbf{Chỉ thị số 21-CT/TW} về tiếp tục đổi mới, phát triển và nâng cao chất lượng giáo dục nghề nghiệp đến năm 2030, tầm nhìn đến năm 2045. Chỉ thị khẳng định phát triển giáo dục nghề nghiệp là nhiệm vụ quan trọng hàng đầu trong chiến lược phát triển nguồn nhân lực chất lượng cao của Đảng, yêu cầu cả hệ thống chính trị phải vào cuộc nhằm nâng cao nhận thức xã hội, đẩy mạnh phân luồng học sinh sau trung học cơ sở và trung học phổ thông, gắn kết chặt chẽ đào tạo với nhu cầu sử dụng của doanh nghiệp.

Trong giai đoạn hiện nay, thực hiện chủ trương của Đảng và Nhà nước về tiếp tục sắp xếp, tinh gọn tổ chức bộ máy của hệ thống chính trị hoạt động hiệu lực, hiệu quả, chức năng quản lý nhà nước về giáo dục nghề nghiệp đang tiếp tục được nghiên cứu, rà soát và điều chỉnh theo hướng thống nhất về mặt thể chế quản lý nhà nước về giáo dục đào tạo, khẳng định sự vận động liên tục của một hệ thống luôn tìm kiếm mô hình tối ưu nhất để phụng sự cho công cuộc phát triển hùng cường của đất nước.
